\documentclass[11pt]{article}
\usepackage{amsmath, amssymb}
\usepackage{geometry}
\usepackage{hyperref}
\geometry{margin=2.5cm}

\title{Derivation: Majorana hybridization splitting $E_0(L) \approx 2t\bigl(|\mu|/2t\bigr)^L$}
\author{Seminar notes}
\date{\today}

\begin{document}
\maketitle

\section{Setup}

We work with the Kitaev chain at the symmetric point $t = \Delta$, open boundary
conditions (OBC), chemical potential $|\mu| < 2t$ (topological phase).
The goal is to derive the finite-size energy of the near-zero Majorana mode.

\section{Edge-mode wavefunction}

\subsection{Bulk zero-energy condition}

A Majorana zero mode satisfies $H\psi = 0$ with a wavefunction that decays
into the bulk from the left edge.  Writing the Ansatz $\psi_j \sim \lambda^j$
and substituting into the bulk BdG equations gives the characteristic equation
\begin{equation}
  \bigl[-\mu - t(\lambda + \lambda^{-1})\bigr]^2
  + \Delta^2(\lambda - \lambda^{-1})^2 = 0.
\end{equation}

\subsection{Solution at the symmetric point $t = \Delta$}

At $t = \Delta$ the equation factors:
\begin{equation}
  \Bigl[-\mu - (t-\Delta)(\lambda+\lambda^{-1}) - (t+\Delta)\lambda^{-1}\cdot\lambda \Bigr]\cdots = 0.
\end{equation}
More directly, substituting $t=\Delta$ one finds two branches:
\begin{equation}
  \lambda_+ = \frac{\mu}{2t}, \qquad \lambda_- = \frac{2t}{\mu}.
\end{equation}
For the \emph{left} edge mode we need $|\lambda| < 1$, so we take
\begin{equation}
  \lambda = \frac{\mu}{2t}, \qquad |\lambda| = \frac{|\mu|}{2t} < 1
  \quad \text{(satisfied inside the topological phase)}.
\end{equation}
The wavefunction is \textbf{purely exponential} (no oscillatory factor at the
symmetric point):
\begin{equation}
  \psi_j^L \propto \lambda^{j-1} = \left(\frac{\mu}{2t}\right)^{j-1},
  \qquad j = 1, 2, \ldots, L.
\end{equation}
The coherence length is
\begin{equation}
  \xi = \frac{1}{\ln(2t/|\mu|)}.
\end{equation}

\subsection{Right edge mode}

By the symmetry $j \leftrightarrow L+1-j$ the right edge mode is
\begin{equation}
  \psi_j^R \propto \lambda^{L-j},
  \qquad j = 1, \ldots, L.
\end{equation}

\section{Finite-size splitting}

\subsection{First-order perturbation theory}

In an infinite chain the two modes $\psi^L$ and $\psi^R$ are degenerate at
$E = 0$.  In a finite chain of length $L$ they overlap, lifting the degeneracy
by the matrix element
\begin{equation}
  \delta E = \frac{|\langle \psi^L | H | \psi^R \rangle|}
                  {\|\psi^L\|\,\|\psi^R\|}.
\end{equation}

\subsection{Why the overlap proxies the matrix element}

The two quantities $\langle \psi^L | H | \psi^R \rangle$ and
$\langle \psi^L | \psi^R \rangle$ are \emph{not} literally equal, but they
carry the same exponential dependence on $L$, which is all that is needed
to determine the decay rate.

\paragraph{Matrix element.}
Both $\psi^L$ and $\psi^R$ are exact zero modes of the \emph{infinite}-chain
Hamiltonian $H_\infty$, so $H_\infty\psi^{L,R}=0$.
The finite-chain Hamiltonian differs from $H_\infty$ only at the boundaries:
open boundary conditions remove the hopping and pairing terms that would
couple to the fictitious sites $j=0$ and $j=L+1$.
Therefore $H|\psi^R\rangle \neq 0$ only near the \emph{left} boundary, where
the right-edge mode has leaked amplitude $\psi^R_{j=1}\sim\lambda^{L-1}$.
The matrix element is dominated by this boundary contribution:
\begin{equation}
  \langle \psi^L | H | \psi^R \rangle
  \;\sim\; \psi^L_1 \cdot (H\psi^R)_1
  \;\sim\; 1 \cdot \lambda^{L-1}
  \;=\; \lambda^{L-1}.
\end{equation}

\paragraph{Overlap integral.}
Summing over all sites gives
\begin{equation}
  \langle \psi^L | \psi^R \rangle
  = \sum_{j=1}^{L} \lambda^{j-1}\,\lambda^{L-j}
  = L\,\lambda^{L-1}.
\end{equation}

\paragraph{Comparison.}
The two expressions differ by a factor of $L$ — polynomial, not exponential.
For large $L$ both vanish as $\lambda^L$, so they encode the same coherence
length $\xi = 1/\ln(2t/|\mu|)$.
The notes use the overlap because it is easier to evaluate; the polynomial
factor $L$ (together with the normalization) is absorbed into the
undetermined prefactor $C$, which is then fixed to $C=2t$ by the exact
sweet-spot calculation.

\subsection{Overlap integral}

The unnormalized overlap is dominated by the ``mid-chain'' region where both
wavefunctions are appreciable:
\begin{equation}
  \langle \psi^L | \psi^R \rangle
  \sim \sum_{j=1}^{L} \lambda^{j-1} \cdot \lambda^{L-j}
  = \lambda^{L-1} \sum_{j=1}^{L} 1
  = L\,\lambda^{L-1}.
\end{equation}

\subsection{Normalisation}

Each mode is normalised over the chain:
\begin{equation}
  \|\psi^L\|^2
  = \sum_{j=1}^{L} \lambda^{2(j-1)}
  = \frac{1 - \lambda^{2L}}{1 - \lambda^2}
  \xrightarrow{L\to\infty} \frac{1}{1-\lambda^2}.
\end{equation}
For $L \gg \xi$ (well inside the topological phase) this is essentially
$\|\psi^L\| \approx 1/\sqrt{1-\lambda^2}$, independent of $L$.

\subsection{Result}

Combining and absorbing the normalisation into the prefactor:
\begin{equation}
  \delta E
  \sim \frac{L\,\lambda^{L-1}}{1/(1-\lambda^2)}
  = L\,(1-\lambda^2)\,\lambda^{L-1}.
\end{equation}
For large $L$ the factor $L\lambda^{L-1} \to 0$ because the exponential decay
wins over the linear growth.  Absorbing the subleading $L$ dependence and
the $(1-\lambda^2)$ prefactor into an overall constant $C$ of order unity:
\begin{equation}
  \delta E \approx C \cdot \lambda^L = C\left(\frac{|\mu|}{2t}\right)^L.
\end{equation}
An exact calculation at the sweet spot $\mu \to 0$, $t = \Delta$ fixes the
prefactor to $C = 2t$, giving the standard result
\begin{equation}
  \boxed{E_0(L) \approx 2t\left(\frac{|\mu|}{2t}\right)^L.}
\end{equation}

\section{Consequences}

\begin{itemize}
  \item \textbf{Exponential decay with $L$:} on a semilog plot $\ln E_0$ vs $L$
    is a straight line with slope $\ln(|\mu|/2t) < 0$.
  \item \textbf{Slope encodes the coherence length:}
    $\text{slope} = -1/\xi = \ln(|\mu|/2t)$.
    Measuring the slope from the numerical data directly gives $\xi$.
  \item \textbf{Sweet spot $\mu = 0$:} $E_0 = 0$ exactly for any $L$
    (the two edge modes have zero overlap because $\lambda = 0$).
  \item \textbf{Phase boundary $|\mu| \to 2t$:} $\lambda \to 1$, so $\xi \to \infty$
    and $E_0 \to 2t$ regardless of $L$ — the zero mode is never protected
    near the transition for any finite chain.
  \item \textbf{Deviations for $t \neq \Delta$:} the wavefunction acquires an
    oscillatory factor $e^{ik_0 j}$ and the splitting becomes
    $E_0 \sim e^{-L/\xi}\cos(k_0 L + \phi)$, leading to
    oscillations around the exponential envelope.  The symmetric point
    $t = \Delta$ is special in that $k_0 = 0$ exactly.
\end{itemize}

\section{Conceptual clarifications}

\paragraph{Are $\psi^L$ and $\psi^R$ two separate states or parts of the
system wavefunction?}
They are two separate single-particle wavefunctions — $\psi^L$ localized at
the left end, $\psi^R$ at the right.
Each is an approximate zero-energy eigenstate of the finite chain on its own.
The true eigenstates of the finite chain are their symmetric and antisymmetric
combinations $(\psi^L \pm \psi^R)/\sqrt{2}$, split by $\pm E_0$.

\paragraph{Can we say $E < 0$ are holes, $E > 0$ are electrons, and $E \approx 0$
are Majoranas?}
Not quite.
In the BdG formalism the eigenstates are \textbf{Bogoliubov quasiparticles}
— superpositions of electrons and holes — at both positive and negative
energies.
The negative-energy branch is not ``holes''; it is the redundant
particle-hole conjugate copy of the positive branch, forced by the doubling
of the Hilbert space in the Nambu spinor construction.
The physical quasiparticle spectrum is the positive side only.
The statement about near-zero modes being Majoranas is correct: a Majorana
mode is its own particle-hole conjugate, which is why it sits at $E = 0$.

\paragraph{If the negative branch is a redundant copy, why is the difference
between $+E$ and $-E$ not zero?}
``Redundant'' means no new \emph{physics}, not the same number.
The negative branch is exactly $-E$ — a mirror image about zero — which is
precisely what particle-hole symmetry guarantees.
Knowing the positive eigenvalues tells you everything; the negative ones add
no new information because they are just the negatives.
The Majorana zero mode is the special case where $+E = -E = 0$, so the
difference actually \emph{is} zero — that is what makes it special.

\paragraph{How does Majorana hybridization splitting fit in?}
In a finite chain the two Majorana modes overlap, so the energy is lifted
from exactly $0$ to $\pm E_0$.
They are still a valid particle-hole conjugate pair (differing by $2E_0$),
but no longer pinned at zero.
They are only true Majorana zero modes in the limit $L \to \infty$ where
$E_0 \to 0$.
The splitting $E_0 \sim (|\mu|/2t)^L$ measures how far the finite system is
from that ideal.

\end{document}
